\begin{pastexam}[label={px:gw3-113-ds-02}]{關務三等資訊處理 113 年 第二題}
（申論；20 分）

一串含有 8 個整數的資料：48, 55, 10, 88, 26, 80, 35, 40，請以第一個數值為樞紐（pivot）進行快速排序法（quick sort），將這些
資料由小排到大，並寫出詳細的比較過程。（20 分）
\end{pastexam}

\begin{pastsoln}
答案：共分割四次，依序得到
$[26, 40, 10, 35]\ \fbox{48}\ [80, 88, 55]$ $\to$ $10\ \fbox{26}\ [40, 35]$ $\to$ $\fbox{40}$ 歸位 $\to$ $\fbox{80}$ 歸位，
最後 $10, 26, 35, 40, 48, 55, 80, 88$。

\textbf{分割方法}（資結書第~14.5~節）：子陣列 $x[\mathit{lb}..\mathit{ub}]$ 以第一個數 $p = x[\mathit{lb}]$ 為樞紐。
\textit{down} 從左往右，跳過 $\le p$ 的數，停在第一個 $> p$ 的數（最多到 $\mathit{ub}$）；\textit{up} 從右往左，跳過 $> p$
的數，停在第一個 $\le p$ 的數。 兩者沒交錯（$\mathit{down} < \mathit{up}$）就交換，繼續找；交錯後把樞紐和 $x[\mathit{up}]$
交換，樞紐就到了最終位置，左邊都 $\le p$、右邊都 $> p$。 接著對左右兩段遞迴。 下面索引從 $0$ 起算，方框是這次的樞紐。

\textbf{第 1 次分割}：$x[0..7] = 48, 55, 10, 88, 26, 80, 35, 40$，$p = 48$。
\begin{itemize}
\item \textit{down}：$48 \le 48$ 跳過，$55 > 48$ 停在 $1$。 \textit{up}：$40 \le 48$ 停在 $7$。 $1 < 7$，交換 $55 \leftrightarrow 40$：
$48, 40, 10, 88, 26, 80, 35, 55$。
\item \textit{down}：$40, 10 \le 48$ 跳過，$88 > 48$ 停在 $3$。 \textit{up}：$55 > 48$ 跳過，$35 \le 48$ 停在 $6$。 $3 < 6$，交換
$88 \leftrightarrow 35$：$48, 40, 10, 35, 26, 80, 88, 55$。
\item \textit{down}：$35, 26 \le 48$ 跳過，$80 > 48$ 停在 $5$。 \textit{up}：$88, 80 > 48$ 跳過，$26 \le 48$ 停在 $4$。 $5 > 4$，
交錯，結束。
\item 樞紐 $48$ 和 $x[4] = 26$ 交換：$26, 40, 10, 35, \fbox{48}, 80, 88, 55$。
\end{itemize}

\textbf{第 2 次分割}：左段 $x[0..3] = 26, 40, 10, 35$，$p = 26$。
\begin{itemize}
\item \textit{down}：$26$ 跳過，$40 > 26$ 停在 $1$。 \textit{up}：$35 > 26$ 跳過，$10 \le 26$ 停在 $2$。 交換 $40 \leftrightarrow 10$：
$26, 10, 40, 35$。
\item \textit{down}：$10$ 跳過，$40 > 26$ 停在 $2$。 \textit{up}：$40 > 26$ 跳過，$10$ 停在 $1$。 交錯，結束。
\item 樞紐 $26$ 和 $x[1] = 10$ 交換：$10, \fbox{26}, 40, 35$。 左段只剩 $10$，不用再排。
\end{itemize}

\textbf{第 3 次分割}：$x[2..3] = 40, 35$，$p = 40$。 \textit{down}：$40$ 跳過，到了右端 $3$ 停下。 \textit{up}：$35 \le 40$ 停在
$3$。 沒有 $\mathit{down} < \mathit{up}$，樞紐和 $x[3] = 35$ 交換：$35, \fbox{40}$。

\textbf{第 4 次分割}：右段 $x[5..7] = 80, 88, 55$，$p = 80$。
\begin{itemize}
\item \textit{down}：$80$ 跳過，$88 > 80$ 停在 $6$。 \textit{up}：$55 \le 80$ 停在 $7$。 交換 $88 \leftrightarrow 55$：$80, 55, 88$。
\item \textit{down}：$55$ 跳過，停在 $7$（$88 > 80$）。 \textit{up}：$88 > 80$ 跳過，$55$ 停在 $6$。 交錯，結束。
\item 樞紐 $80$ 和 $x[6] = 55$ 交換：$55, \fbox{80}, 88$。 兩邊各剩一個數。
\end{itemize}

\textbf{整理}：
\begin{center}
\begin{tabular}{cll}
\toprule
次 & 範圍與樞紐 & 分割後的整個陣列 \\
\midrule
原始 & & 48, 55, 10, 88, 26, 80, 35, 40 \\
1 & $[0..7]$，48 & 26, 40, 10, 35, \fbox{48}, 80, 88, 55 \\
2 & $[0..3]$，26 & 10, \fbox{26}, 40, 35, 48, 80, 88, 55 \\
3 & $[2..3]$，40 & 10, 26, 35, \fbox{40}, 48, 80, 88, 55 \\
4 & $[5..7]$，80 & 10, 26, 35, 40, 48, 55, \fbox{80}, 88 \\
\bottomrule
\end{tabular}
\end{center}
排序結果 $10, 26, 35, 40, 48, 55, 80, 88$。 每次分割把樞紐放到最終位置；這組資料的樞紐 $48$ 剛好把資料分成 $4$ 和 $3$ 個，
接近平分，遞迴只有三層。 （其他分割寫法，例如 Lomuto，中間過程會不同，最後結果一樣。）
\end{pastsoln}
